% !TeX root = ../tfg.tex
% !TeX encoding = utf8

\chapter{Detección OOD con ConvNeXt-S}\label{ap:convnext}

% Fuentes: OpenOOD/results_ood_convnext (resumen.md, decisiones.md y CSV). Tablas generadas por
% scripts/generar_tablas.py (tD_*).
% Puntos de la introducción:
% - Se repite la tabla principal del \cref{cap:deteccion-ood} con ConvNeXt-S, con los mismos conjuntos,
%   detectores, referencias y bootstrap, para comprobar qué conclusiones dependen del modelo.
% - No se repiten la fuga por lámina, la predicción conforme, Mondrian, el anexo NCT ni las corrupciones.
% - Las cabezas no son equivalentes: en el Swin, las 768 características salen de la red troncal y entran
%   directamente en la única capa lineal; en ConvNeXt, la cabeza de fastai tiene dos capas lineales.
% - Espacio principal B: salida de la ReLU, 512 dimensiones, entrada de la última lineal, con la BatchNorm
%   plegada (exacta en evaluación). Es lo que exigen ViM y ReAct, y coincide con la hipótesis post-ReLU de
%   ReAct, ASH y DICE.
% - ViM con dimensión 171, la misma proporción 1/3 que 256/768 en el Swin.
% - ConvNeXt evaluado en CPU (fp32) y el Swin en GPU. Exactitud de test 0,96695: 2 imágenes más que en la
%   fase 1 (0,9668).
% - Cabeza de fastai y entrada de 256 px: \cref{sec:protocolo}.
\pendiente{introducción del autor}

\section{Configuración}\label{sec:apd-config}

% Generado por scripts/generar_tablas.py a partir de results_ood_convnext/configuracion_convnext.csv. No editar a mano.
\begin{table}[tbp]
\centering
\footnotesize
\caption{Configuración de la evaluación OOD con ConvNeXt-S y diferencias con la del Swin-S. B y C son los espacios de características; $W'$ y $b'$, la última capa lineal con la BatchNorm plegada.}
\label{tab:convnext-config}
\begin{tabularx}{\textwidth}{@{}l>{\raggedright\arraybackslash}X@{}}
\toprule
Parámetro & Valor \\
\midrule
Modelo & ConvNeXt-S (fastai); mismo entrenamiento que en el \cref{sec:protocolo} \\
Cabeza & fastai: AdaptiveConcatPool2d (1.536) $\to$ BN1d $\to$ Dropout(0,25) $\to$ Linear(1.536, 512) $\to$ ReLU $\to$ BN1d(512) $\to$ Dropout(0,5) $\to$ Linear(512, 14). Swin: Linear(768, 14) de timm sobre la red troncal \\
Espacio principal & B, 512: salida de la ReLU, BN1d(512) plegada en W$'$, b$'$. Máx.\ $|W'B+b'-\mathrm{logits}|$ en los 11 conjuntos: $4{,}88\cdot 10^{-6}$ \\
Espacio de sensibilidad & C, 1.536: concat pooling, antes de BN1d(1.536) \\
Dimensión de ViM & 171 (sensibilidad: 256; Swin: 256 de 768) \\
$K$ de \stoodx & 500 (coseno, clase predicha, 100 \% de características) \\
Ejecución & nodo zeus, CPU (Intel Xeon Silver 4114, 2,20 GHz), 12 hilos, fp32 (Swin en GPU) \\
Tiempo total & 8 h 56 min 12 s (extracción 7,85 h) \\
Memoria (MaxRSS) & 6,2 GiB (6.480.436 KB) \\
Exactitud de test & 0,96695 (12.201/12.618; fase 1: 0,9668, +2 imágenes) \\
\bottomrule
\end{tabularx}
\end{table}


% Cifras clave (configuracion_convnext.csv; comprobacion_hw.json; config_detectores.json):
% - máx. |W'B + b' - logits| = 4,88e-6 (id_train); < 1e-4 en los 11 conjuntos.
% - Exactitud de test 0,96695 (12.201/12.618); fase 1: 0,9668.
% - Exactitud por detector: igual salvo ReAct 0,9667, ASH 0,9649 y DICE 0,2845.
% - zeus, CPU, 12 hilos, fp32; 8 h 56 min 12 s (extracción 7,85 h); MaxRSS 6,2 GiB.
\pendiente{comentario breve}

\section{Resultados}\label{sec:apd-resultados}

% Generado por scripts/generar_tablas.py a partir de results_ood_convnext/tabla_principal_convnext.csv. No editar a mano.
\begin{table}[tbp]
\centering
\scriptsize
\setlength{\tabcolsep}{2pt}
\caption{ConvNeXt-S, espacio B: \AUROC (más alto es mejor) de los 14 métodos en los nueve conjuntos. Clase positiva: OOD. En negrita, el mejor valor de cada conjunto. IC del 95\,\% por \emph{bootstrap} (1.000 remuestreos estratificados, semilla 0). Los IC del resto de conjuntos están en \texttt{tabla\_principal\_convnext.csv}. $^{a}$Modifican la predicción: la exactitud ID baja al 96{,}5\,\% (ASH) y al 28{,}5\,\% (DICE, con el postprocesador de OpenOOD sin cambios).}
\label{tab:convnext-auroc}
\begin{tabular}{lccccccccc}
\toprule
Método & TUM+NORM [IC 95\,\%] & TUM & NORM & STR & LYM & MUS & PANDA & Control 128\,px & DTD \\
\midrule
\multicolumn{10}{@{}l}{\textit{Basados en logits}} \\
MSP & 0{,}954 [0{,}951; 0{,}956] & 0{,}940 & 0{,}977 & 0{,}948 & \textbf{0{,}974} & 0{,}891 & 0{,}926 & 0{,}660 & 0{,}988 \\
MaxLogit & 0{,}867 [0{,}864; 0{,}871] & 0{,}877 & 0{,}851 & 0{,}782 & 0{,}893 & 0{,}723 & 0{,}689 & 0{,}561 & 0{,}918 \\
Energy & 0{,}833 [0{,}829; 0{,}837] & 0{,}846 & 0{,}812 & 0{,}740 & 0{,}841 & 0{,}691 & 0{,}656 & 0{,}554 & 0{,}865 \\
\midrule
\multicolumn{10}{@{}l}{\textit{Modificación de activaciones}} \\
ReAct & 0{,}830 [0{,}826; 0{,}834] & 0{,}843 & 0{,}809 & 0{,}729 & 0{,}838 & 0{,}681 & 0{,}659 & 0{,}553 & 0{,}860 \\
ASH$^{a}$ & 0{,}828 [0{,}823; 0{,}832] & 0{,}829 & 0{,}826 & 0{,}754 & 0{,}723 & 0{,}628 & 0{,}447 & 0{,}516 & 0{,}899 \\
DICE$^{a}$ & 0{,}616 [0{,}610; 0{,}622] & 0{,}594 & 0{,}651 & 0{,}533 & 0{,}209 & 0{,}420 & 0{,}211 & 0{,}501 & 0{,}610 \\
\midrule
\multicolumn{10}{@{}l}{\textit{Características y distancias}} \\
KNN & \textbf{0{,}972} [0{,}971; 0{,}974] & \textbf{0{,}963} & \textbf{0{,}987} & 0{,}960 & 0{,}969 & 0{,}898 & 0{,}885 & 0{,}717 & \textbf{0{,}997} \\
MDS & 0{,}953 [0{,}950; 0{,}955] & 0{,}952 & 0{,}954 & 0{,}948 & 0{,}968 & 0{,}898 & 0{,}886 & 0{,}709 & 0{,}981 \\
RMDS & 0{,}958 [0{,}955; 0{,}960] & 0{,}955 & 0{,}963 & 0{,}940 & 0{,}970 & 0{,}898 & 0{,}921 & 0{,}714 & 0{,}969 \\
ViM & 0{,}939 [0{,}936; 0{,}942] & 0{,}935 & 0{,}945 & 0{,}930 & 0{,}952 & 0{,}876 & 0{,}833 & 0{,}677 & 0{,}977 \\
SHE (prod.\ escalar) & 0{,}725 [0{,}720; 0{,}730] & 0{,}689 & 0{,}783 & 0{,}775 & 0{,}628 & 0{,}756 & 0{,}666 & 0{,}543 & 0{,}888 \\
SHE (coseno) & 0{,}942 [0{,}940; 0{,}944] & 0{,}922 & 0{,}976 & 0{,}931 & 0{,}940 & 0{,}845 & 0{,}793 & 0{,}627 & 0{,}991 \\
SHE (euclídea) & 0{,}939 [0{,}936; 0{,}941] & 0{,}928 & 0{,}957 & 0{,}919 & 0{,}957 & 0{,}832 & 0{,}798 & 0{,}620 & 0{,}977 \\
\midrule
\multicolumn{10}{@{}l}{\textit{STOOD-X}} \\
STOOD-X (100\,\%) & 0{,}949 [0{,}947; 0{,}951] & 0{,}930 & 0{,}981 & \textbf{0{,}986} & 0{,}957 & \textbf{0{,}952} & \textbf{0{,}975} & \textbf{0{,}782} & 0{,}993 \\
\bottomrule
\end{tabular}
\end{table}


% Generado por scripts/generar_tablas.py a partir de results_ood_convnext/tabla_principal_convnext.csv. No editar a mano.
\begin{table}[tbp]
\centering
\scriptsize
\setlength{\tabcolsep}{2.5pt}
\caption{ConvNeXt-S, espacio B: \FPRn{95} (más bajo es mejor) de los 14 métodos en los nueve conjuntos y \CCRn{5} (más alto es mejor) en TUM+NORM. Clase positiva: OOD. En negrita, el mejor valor de cada columna. Los IC de \FPRn{95} están en \texttt{tabla\_principal\_convnext.csv}. $^{a}$Modifican la predicción: la exactitud ID baja al 96{,}5\,\% (ASH) y al 28{,}5\,\% (DICE, con el postprocesador de OpenOOD sin cambios).}
\label{tab:convnext-fpr}
\begin{tabular}{lcccccccccc}
\toprule
& \multicolumn{9}{c}{\FPRn{95} $\downarrow$} & \CCRn{5} $\uparrow$ \\
\cmidrule(lr){2-10}\cmidrule(lr){11-11}
Método & TUM+NORM & TUM & NORM & STR & LYM & MUS & PANDA & Control 128\,px & DTD & TUM+NORM \\
\midrule
\multicolumn{11}{@{}l}{\textit{Basados en logits}} \\
MSP & 0{,}108 & 0{,}170 & 0{,}084 & 0{,}160 & \textbf{0{,}086} & 0{,}479 & 0{,}176 & 0{,}910 & 0{,}049 & 0{,}884 \\
MaxLogit & 0{,}584 & 0{,}524 & 0{,}751 & 0{,}716 & 0{,}451 & 0{,}610 & 0{,}546 & 0{,}961 & 0{,}357 & 0{,}413 \\
Energy & 0{,}660 & 0{,}560 & 0{,}772 & 0{,}753 & 0{,}493 & 0{,}645 & 0{,}571 & 0{,}961 & 0{,}376 & 0{,}338 \\
\midrule
\multicolumn{11}{@{}l}{\textit{Modificación de activaciones}} \\
ReAct & 0{,}692 & 0{,}569 & 0{,}766 & 0{,}760 & 0{,}484 & 0{,}721 & 0{,}562 & 0{,}960 & 0{,}378 & 0{,}306 \\
ASH$^{a}$ & 0{,}715 & 0{,}713 & 0{,}718 & 0{,}668 & 0{,}922 & 0{,}922 & 0{,}972 & 0{,}969 & 0{,}379 & 0{,}281 \\
DICE$^{a}$ & 0{,}766 & 0{,}762 & 0{,}770 & 0{,}842 & 0{,}990 & 0{,}962 & 0{,}948 & 0{,}944 & 0{,}763 & 0{,}005 \\
\midrule
\multicolumn{11}{@{}l}{\textit{Características y distancias}} \\
KNN & 0{,}136 & 0{,}193 & \textbf{0{,}051} & 0{,}200 & 0{,}127 & 0{,}504 & 0{,}364 & \textbf{0{,}692} & \textbf{0{,}013} & 0{,}859 \\
MDS & 0{,}143 & 0{,}159 & 0{,}118 & 0{,}168 & 0{,}119 & 0{,}363 & 0{,}303 & 0{,}734 & 0{,}047 & 0{,}851 \\
RMDS & \textbf{0{,}103} & \textbf{0{,}107} & 0{,}098 & 0{,}156 & 0{,}097 & 0{,}311 & 0{,}191 & 0{,}843 & 0{,}069 & \textbf{0{,}890} \\
ViM & 0{,}193 & 0{,}209 & 0{,}162 & 0{,}207 & 0{,}161 & 0{,}400 & 0{,}411 & 0{,}802 & 0{,}066 & 0{,}803 \\
SHE (prod.\ escalar) & 0{,}892 & 0{,}941 & 0{,}710 & 0{,}664 & 0{,}911 & 0{,}646 & 0{,}587 & 0{,}926 & 0{,}408 & 0{,}103 \\
SHE (coseno) & 0{,}336 & 0{,}429 & 0{,}098 & 0{,}401 & 0{,}299 & 0{,}757 & 0{,}565 & 0{,}853 & 0{,}031 & 0{,}659 \\
SHE (euclídea) & 0{,}251 & 0{,}323 & 0{,}157 & 0{,}348 & 0{,}184 & 0{,}761 & 0{,}625 & 0{,}857 & 0{,}075 & 0{,}746 \\
\midrule
\multicolumn{11}{@{}l}{\textit{STOOD-X}} \\
STOOD-X (100\,\%) & 0{,}242 & 0{,}421 & 0{,}158 & \textbf{0{,}053} & 0{,}236 & \textbf{0{,}232} & \textbf{0{,}078} & 0{,}729 & 0{,}025 & 0{,}752 \\
\bottomrule
\end{tabular}
\end{table}


% Cifras clave (tabla_principal_convnext.csv):
% - TUM+NORM, AUROC: KNN 0,972 [0,971; 0,974] (1.º), RMDS 0,958, MSP 0,954, MDS 0,953, STOOD-X 0,949;
%   ViM 0,939 (8.º).
% - Mejor AUROC por conjunto: KNN en TUM+NORM, TUM (0,963), NORM (0,987) y DTD (0,997); STOOD-X en STR
%   (0,986), MUS (0,952), PANDA (0,975) y control (0,782); MSP en LYM (0,974).
% - TUM+NORM, FPR@95: RMDS 0,103 [0,097; 0,108], MSP 0,108 [0,102; 0,116], KNN 0,136.
% - TUM+NORM, CCR@FPR5: RMDS 0,890, MSP 0,884, KNN 0,859.
% - DICE colapsa: AUROC 0,209 en LYM y 0,211 en PANDA (exactitud ID 0,2845).
\pendiente{comentario breve}

\section{Comparación con el Swin-S}\label{sec:apd-comparacion}

% Generado por scripts/generar_tablas.py a partir de results_ood_convnext/diferencias_auroc.csv. No editar a mano.
\begin{table}[tbp]
\centering
\footnotesize
\setlength{\tabcolsep}{2.5pt}
\caption{Diferencia de \AUROC ConvNeXt-S $-$ Swin-S (positiva: mejor ConvNeXt-S) en los métodos principales, con su IC debajo. IC del 95\,\% por \emph{bootstrap} pareado por imagen (1.000 remuestreos estratificados ID/OOD, semilla 0). Los demás métodos están en \texttt{diferencias\_auroc.csv}.}
\label{tab:convnext-diferencias}
\begin{tabular}{lcccccc}
\toprule
Conjunto & MSP & KNN & MDS & RMDS & ViM & STOOD-X (100\,\%) \\
\midrule
TUM+NORM & +0{,}078 & +0{,}016 & $-$0{,}015 & +0{,}004 & $-$0{,}031 & $-$0{,}004 \\
& {\scriptsize [+0{,}074; +0{,}082]} & {\scriptsize [+0{,}014; +0{,}018]} & {\scriptsize [$-$0{,}017; $-$0{,}013]} & {\scriptsize [+0{,}002; +0{,}007]} & {\scriptsize [$-$0{,}033; $-$0{,}028]} & {\scriptsize [$-$0{,}006; $-$0{,}002]} \\[2pt]
TUM & +0{,}066 & +0{,}011 & $-$0{,}012 & +0{,}007 & $-$0{,}031 & $-$0{,}001 \\
& {\scriptsize [+0{,}061; +0{,}072]} & {\scriptsize [+0{,}009; +0{,}013]} & {\scriptsize [$-$0{,}014; $-$0{,}010]} & {\scriptsize [+0{,}004; +0{,}011]} & {\scriptsize [$-$0{,}034; $-$0{,}029]} & {\scriptsize [$-$0{,}004; +0{,}002]} \\[2pt]
NORM & +0{,}097 & +0{,}024 & $-$0{,}019 & $-$0{,}001 & $-$0{,}030 & $-$0{,}009 \\
& {\scriptsize [+0{,}091; +0{,}103]} & {\scriptsize [+0{,}022; +0{,}026]} & {\scriptsize [$-$0{,}022; $-$0{,}017]} & {\scriptsize [$-$0{,}004; +0{,}001]} & {\scriptsize [$-$0{,}033; $-$0{,}028]} & {\scriptsize [$-$0{,}011; $-$0{,}008]} \\[2pt]
STR & +0{,}126 & +0{,}029 & +0{,}005 & +0{,}005 & $-$0{,}015 & +0{,}008 \\
& {\scriptsize [+0{,}119; +0{,}132]} & {\scriptsize [+0{,}026; +0{,}031]} & {\scriptsize [+0{,}002; +0{,}007]} & {\scriptsize [+0{,}002; +0{,}008]} & {\scriptsize [$-$0{,}018; $-$0{,}013]} & {\scriptsize [+0{,}006; +0{,}010]} \\[2pt]
LYM & +0{,}120 & +0{,}021 & +0{,}001 & +0{,}147 & $-$0{,}018 & +0{,}054 \\
& {\scriptsize [+0{,}115; +0{,}125]} & {\scriptsize [+0{,}019; +0{,}024]} & {\scriptsize [$-$0{,}001; +0{,}003]} & {\scriptsize [+0{,}141; +0{,}153]} & {\scriptsize [$-$0{,}020; $-$0{,}016]} & {\scriptsize [+0{,}051; +0{,}057]} \\[2pt]
MUS & +0{,}150 & $-$0{,}006 & $-$0{,}015 & $-$0{,}013 & $-$0{,}039 & $-$0{,}015 \\
& {\scriptsize [+0{,}144; +0{,}157]} & {\scriptsize [$-$0{,}010; $-$0{,}003]} & {\scriptsize [$-$0{,}017; $-$0{,}012]} & {\scriptsize [$-$0{,}017; $-$0{,}010]} & {\scriptsize [$-$0{,}042; $-$0{,}036]} & {\scriptsize [$-$0{,}018; $-$0{,}013]} \\[2pt]
PANDA & +0{,}083 & $-$0{,}007 & $-$0{,}025 & +0{,}009 & $-$0{,}079 & +0{,}079 \\
& {\scriptsize [+0{,}077; +0{,}088]} & {\scriptsize [$-$0{,}010; $-$0{,}003]} & {\scriptsize [$-$0{,}028; $-$0{,}022]} & {\scriptsize [+0{,}005; +0{,}013]} & {\scriptsize [$-$0{,}083; $-$0{,}075]} & {\scriptsize [+0{,}075; +0{,}084]} \\[2pt]
Control 128\,px & +0{,}003 & $-$0{,}042 & $-$0{,}070 & $-$0{,}039 & $-$0{,}112 & $-$0{,}112 \\
& {\scriptsize [$-$0{,}006; +0{,}011]} & {\scriptsize [$-$0{,}047; $-$0{,}037]} & {\scriptsize [$-$0{,}074; $-$0{,}066]} & {\scriptsize [$-$0{,}045; $-$0{,}034]} & {\scriptsize [$-$0{,}116; $-$0{,}108]} & {\scriptsize [$-$0{,}117; $-$0{,}107]} \\[2pt]
DTD & +0{,}112 & +0{,}060 & +0{,}025 & +0{,}027 & +0{,}018 & +0{,}002 \\
& {\scriptsize [+0{,}106; +0{,}118]} & {\scriptsize [+0{,}057; +0{,}064]} & {\scriptsize [+0{,}023; +0{,}027]} & {\scriptsize [+0{,}024; +0{,}030]} & {\scriptsize [+0{,}016; +0{,}021]} & {\scriptsize [+0{,}001; +0{,}004]} \\[2pt]
\bottomrule
\end{tabular}
\end{table}


% Generado por scripts/generar_tablas.py a partir de results_ood_convnext/conclusiones_swin.csv. No editar a mano.
\begin{table}[tbp]
\centering
\scriptsize
\setlength{\tabcolsep}{3pt}
\caption{Conclusiones del \cref{cap:deteccion-ood} con Swin-S, una por fila, y cifras que las sostienen con cada modelo (\AUROC salvo que se indique otra métrica). No concluyente: la ordenación se mantiene, pero los IC se solapan.}
\label{tab:convnext-conclusiones}
\begin{tabularx}{\textwidth}{@{}l>{\raggedright\arraybackslash}p{2.4cm}>{\raggedright\arraybackslash}X>{\raggedright\arraybackslash}X>{\raggedright\arraybackslash}p{1.4cm}@{}}
\toprule
\# & Conclusión con Swin-S & Swin-S & ConvNeXt-S (B) & Estado \\
\midrule
1 & ViM y MDS son los mejores en TUM+NORM & ViM 0,969 (1.º); MDS 0,968 (2.º) & KNN 0,972 (1.º); RMDS 0,958; MSP 0,954; MDS 0,953 (4.º); ViM 0,939 (8.º) & Cambia \\
2a & RMDS tiene el mejor FPR@95 en TUM+NORM & RMDS 0,128 (1.º); ViM 0,143 & RMDS 0,103 [0,097; 0,108] (1.º); MSP 0,108 [0,102; 0,116]; KNN 0,136 & No concluyente \\
2b & RMDS falla en LYM & AUROC 0,824; FPR@95 0,992 & AUROC 0,970; FPR@95 0,097 & Cambia \\
3a & STOOD-X destaca en STR y MUS & 1.º en los dos: STR 0,978; MUS 0,968 & 1.º en los dos: STR 0,986; MUS 0,952 & Se mantiene \\
3b & STOOD-X destaca en NORM y DTD & 1.º en los dos: NORM 0,990; DTD 0,991 & 2.º tras KNN: NORM 0,981 (KNN 0,987); DTD 0,993 (KNN 0,997); IC sin solapamiento & Cambia \\
4 & MSP $>$ MaxLogit $>$ Energy & se cumple en 9/9 conjuntos & se cumple en 9/9 (TUM+NORM: 0,954 $>$ 0,867 $>$ 0,833) & Se mantiene \\
5 & ReAct, ASH y DICE no funcionan (por debajo de MSP) & por debajo de MSP en 9/9 los tres; por debajo de Energy: ReAct 1/9, ASH 9/9, DICE 8/9 & por debajo de MSP en 9/9 los tres; por debajo de Energy: ReAct 8/9, ASH 6/9, DICE 9/9; TUM+NORM: ReAct 0,830, ASH 0,828, DICE 0,616 & Se mantiene \\
6 & LYM $\approx$ TUM/NORM & AUROC(LYM) $-$ AUROC(TUM+NORM) entre $-$0,024 y 0,000 en KNN, MDS, ViM, SHE-cos y SHE-euc; MSP $-$0,022; RMDS $-$0,130; STOOD-X $-$0,050 & entre $-$0,003 y +0,021 en los 8 (KNN $-$0,003; MDS +0,015; RMDS +0,013; ViM +0,013; SHE-cos $-$0,003; SHE-euc +0,018; STOOD-X +0,008; MSP +0,021) & Se mantiene \\
7 & Los métodos de distancia rinden más en PANDA que en el control & AUROC(PANDA) $-$ AUROC(control) de +0,123 a +0,205 en los 6 de distancia; STOOD-X +0,001 & de +0,156 a +0,207 en los 7 (STOOD-X +0,193) & Se mantiene \\
8 & STOOD-X no distingue PANDA del control & AUROC(control $-$, PANDA +) 0,475 [0,467; 0,482] & 0,859 [0,855; 0,864] & Cambia \\
9 & MDS supera a KNN & MDS $-$ KNN $>$ 0 en 9/9; TUM+NORM +0,011 [+0,011; +0,012]; de +0,009 (MUS) a +0,020 (control) & MDS $-$ KNN $<$ 0 con IC sin el 0 en 7/9 (TUM+NORM $-$0,020 [$-$0,022; $-$0,018]; NORM $-$0,033; DTD $-$0,016; LYM $-$0,002 [$-$0,003; $-$0,000]); IC con el 0 en MUS +0,000 [$-$0,002; +0,002] y PANDA +0,001 [$-$0,001; +0,003] & Cambia \\
\bottomrule
\end{tabularx}
\end{table}


% Generado por scripts/generar_tablas.py a partir de results_ood_convnext/espacio_principal/kendall.csv. No editar a mano.
\begin{table}[tbp]
\centering
\footnotesize
\setlength{\tabcolsep}{5pt}
\caption{$\tau$ de Kendall entre las ordenaciones de los métodos por \AUROC con Swin-S y con ConvNeXt-S, en cada conjunto ($\tau=1$: misma ordenación); ConvNeXt-S en el espacio M (ViM, ReAct, ASH y DICE en B). Valores $p$ sin corregir por comparaciones múltiples.}
\label{tab:convnext-kendall}
\begin{tabular}{lcccc}
\toprule
& \multicolumn{2}{c}{14 métodos} & \multicolumn{2}{c}{Sin ASH ni DICE (12)} \\
\cmidrule(lr){2-3}\cmidrule(lr){4-5}
Conjunto & $\tau$ & $p$ & $\tau$ & $p$ \\
\midrule
TUM+NORM & 0{,}670 & 0{,}0005 & 0{,}606 & 0{,}0054 \\
TUM & 0{,}604 & 0{,}0020 & 0{,}515 & 0{,}0210 \\
NORM & 0{,}626 & 0{,}0012 & 0{,}545 & 0{,}0138 \\
STR & 0{,}692 & 0{,}0003 & 0{,}667 & 0{,}0018 \\
LYM & 0{,}582 & 0{,}0030 & 0{,}606 & 0{,}0054 \\
MUS & 0{,}780 & $<0{,}0001$ & 0{,}727 & 0{,}0005 \\
PANDA & 0{,}473 & 0{,}0193 & 0{,}333 & 0{,}1526 \\
Control 128\,px & 0{,}802 & $<0{,}0001$ & 0{,}758 & 0{,}0002 \\
DTD & 0{,}648 & 0{,}0008 & 0{,}606 & 0{,}0054 \\
\bottomrule
\end{tabular}
\end{table}


% Cifras clave (diferencias_auroc.csv, conclusiones_swin.csv, kendall.csv, mds_menos_knn.csv):
% - ConvNeXt - Swin en TUM+NORM: MSP +0,078 [+0,074; +0,082]; KNN +0,016; MDS -0,015; ViM -0,031;
%   RMDS +0,004; STOOD-X -0,004. RMDS en LYM: +0,147 [+0,141; +0,153].
% - MDS - KNN en TUM+NORM: Swin +0,011 [+0,011; +0,012]; ConvNeXt -0,020 [-0,022; -0,018].
% - Kendall (14 métodos): de 0,385 (LYM, p = 0,0617) a 0,846 (control, p < 0,0001); TUM+NORM 0,604
%   (p = 0,0020). Sin ASH ni DICE: NORM 0,303 (p = 0,1970) y LYM 0,333 (p = 0,1526).
% - Conclusiones: se mantienen 4, 5, 6, 7 y 3a (STR, MUS); cambian 1, 2b, 3b (NORM, DTD), 8 y 9;
%   2a, no concluyente.
\pendiente{comentario breve}

\section{\texorpdfstring{\stoodx}{STOOD-X} y la saturación}\label{sec:apd-stoodx}

% Generado por scripts/generar_tablas.py a partir de results_ood_convnext/panda_vs_control.csv y saturacion_stoodx.csv. No editar a mano.
\begin{table}[tbp]
\centering
\footnotesize
\setlength{\tabcolsep}{4pt}
\caption{\stoodx y la saturación con los dos modelos. (a) \AUROC con el control de 128\,px como clase negativa y PANDA como positiva: por encima de 0,5, la puntuación coloca PANDA como más OOD que el control. IC del 95\,\% por \emph{bootstrap} (1.000 remuestreos estratificados, semilla 0). (b) Porcentaje de imágenes con la puntuación $-z$ de \stoodx ($K=500$) saturada ($U=n_1n_2$); en C solo se calculó en el test ID, TUM y NORM.}
\label{tab:convnext-stoodx}

\textit{(a) PANDA frente al control}\par\smallskip
\begin{tabular}{lcc}
\toprule
Método & Swin-S & ConvNeXt-S \\
\midrule
\multicolumn{3}{@{}l}{\textit{Basados en logits}} \\
MSP & 0{,}732 [0{,}725; 0{,}738] & 0{,}803 [0{,}797; 0{,}808] \\
MaxLogit & 0{,}617 [0{,}610; 0{,}625] & 0{,}603 [0{,}596; 0{,}610] \\
Energy & 0{,}600 [0{,}593; 0{,}608] & 0{,}573 [0{,}566; 0{,}581] \\
\midrule
\multicolumn{3}{@{}l}{\textit{Modificación de activaciones}} \\
ReAct & 0{,}639 [0{,}632; 0{,}645] & 0{,}577 [0{,}570; 0{,}585] \\
ASH$^{a}$ & 0{,}339 [0{,}331; 0{,}346] & 0{,}435 [0{,}427; 0{,}443] \\
DICE$^{a}$ & 0{,}580 [0{,}573; 0{,}587] & 0{,}195 [0{,}189; 0{,}201] \\
\midrule
\multicolumn{3}{@{}l}{\textit{Características y distancias}} \\
KNN & 0{,}731 [0{,}725; 0{,}738] & 0{,}736 [0{,}729; 0{,}742] \\
MDS & 0{,}729 [0{,}722; 0{,}735] & 0{,}749 [0{,}742; 0{,}754] \\
RMDS & 0{,}768 [0{,}762; 0{,}774] & 0{,}782 [0{,}776; 0{,}788] \\
ViM & 0{,}721 [0{,}714; 0{,}727] & 0{,}673 [0{,}667; 0{,}680] \\
SHE (prod.\ escalar) & 0{,}424 [0{,}417; 0{,}431] & 0{,}611 [0{,}604; 0{,}619] \\
SHE (coseno) & 0{,}706 [0{,}699; 0{,}712] & 0{,}686 [0{,}679; 0{,}693] \\
SHE (euclídea) & 0{,}726 [0{,}719; 0{,}732] & 0{,}688 [0{,}681; 0{,}695] \\
\midrule
\multicolumn{3}{@{}l}{\textit{STOOD-X}} \\
STOOD-X (100\,\%) & 0{,}475 [0{,}467; 0{,}482] & 0{,}859 [0{,}855; 0{,}864] \\
\bottomrule
\end{tabular}

\bigskip
\textit{(b) Imágenes saturadas}\par\smallskip
\begin{tabular}{lccc}
\toprule
Conjunto & Swin-S & ConvNeXt-S (B) & ConvNeXt-S (C) \\
\midrule
Test ID & 0{,}6\,\% & 0{,}1\,\% & 0{,}0\,\% \\
TUM+NORM & 66{,}3\,\% & 26{,}0\,\% & 9{,}9\,\% \\
TUM & 59{,}0\,\% & 30{,}1\,\% & 9{,}7\,\% \\
NORM & 78{,}2\,\% & 19{,}4\,\% & 10{,}3\,\% \\
STR & 52{,}0\,\% & 24{,}7\,\% & -- \\
LYM & 32{,}1\,\% & 18{,}0\,\% & -- \\
MUS & 44{,}1\,\% & 13{,}2\,\% & -- \\
PANDA & 19{,}4\,\% & 7{,}5\,\% & -- \\
Control 128\,px & 19{,}0\,\% & 1{,}5\,\% & -- \\
DTD & 76{,}3\,\% & 31{,}4\,\% & -- \\
\bottomrule
\end{tabular}
\end{table}


% Cifras clave (panda_vs_control.csv, saturacion_stoodx.csv):
% - PANDA frente al control, STOOD-X: Swin 0,475 [0,467; 0,482]; ConvNeXt 0,859 [0,855; 0,864].
% - Saturadas: TUM+NORM 66,3 % (Swin) -> 26,0 % (B) -> 9,9 % (C); control 19,0 % -> 1,5 %;
%   DTD 76,3 % -> 31,4 %; test ID 0,6 % -> 0,1 %.
% - El -z mínimo teórico (-38,6911 con K = 500) se alcanza en todos los conjuntos de los dos modelos.
% - En ConvNeXt no se repitió el análisis de empates (en el Swin acotó el efecto en el AUROC a <= 0,0031).
\pendiente{comentario breve}

\section{Sensibilidad}\label{sec:apd-sensibilidad}

% Generado por scripts/generar_tablas.py a partir de results_ood_convnext/sensibilidad_pareada.csv. No editar a mano.
\begin{table}[tbp]
\centering
\footnotesize
\setlength{\tabcolsep}{4pt}
\caption{Sensibilidad de ConvNeXt-S: diferencia de \AUROC entre el espacio C (1.536 dimensiones) y el B (512), y entre ViM con 256 y con 171 dimensiones en B (positiva: mejor C o ViM 256), con su IC debajo. IC del 95\,\% por \emph{bootstrap} pareado por imagen (1.000 remuestreos estratificados ID/OOD, semilla 0). \stoodx en C solo se calculó en TUM y NORM.}
\label{tab:convnext-sensibilidad}
\begin{tabular}{lccccc}
\toprule
& \multicolumn{4}{c}{C $-$ B} & \\
\cmidrule(lr){2-5}
Conjunto & KNN & MDS & RMDS & STOOD-X & ViM 256 $-$ 171 \\
\midrule
TUM+NORM & $-$0{,}016 & $-$0{,}019 & +0{,}005 & $-$0{,}041 & +0{,}002 \\
& {\scriptsize [$-$0{,}018; $-$0{,}014]} & {\scriptsize [$-$0{,}021; $-$0{,}017]} & {\scriptsize [+0{,}003; +0{,}006]} & {\scriptsize [$-$0{,}043; $-$0{,}039]} & {\scriptsize [+0{,}002; +0{,}002]} \\[2pt]
TUM & $-$0{,}018 & $-$0{,}034 & +0{,}001 & $-$0{,}057 & +0{,}004 \\
& {\scriptsize [$-$0{,}020; $-$0{,}016]} & {\scriptsize [$-$0{,}036; $-$0{,}032]} & {\scriptsize [$-$0{,}001; +0{,}003]} & {\scriptsize [$-$0{,}060; $-$0{,}055]} & {\scriptsize [+0{,}003; +0{,}004]} \\[2pt]
NORM & $-$0{,}013 & +0{,}005 & +0{,}011 & $-$0{,}014 & $-$0{,}001 \\
& {\scriptsize [$-$0{,}014; $-$0{,}011]} & {\scriptsize [+0{,}003; +0{,}008]} & {\scriptsize [+0{,}009; +0{,}012]} & {\scriptsize [$-$0{,}016; $-$0{,}012]} & {\scriptsize [$-$0{,}001; $-$0{,}001]} \\[2pt]
STR & +0{,}004 & +0{,}006 & +0{,}016 & -- & +0{,}004 \\
& {\scriptsize [+0{,}002; +0{,}005]} & {\scriptsize [+0{,}004; +0{,}008]} & {\scriptsize [+0{,}015; +0{,}018]} &  & {\scriptsize [+0{,}004; +0{,}005]} \\[2pt]
LYM & $-$0{,}003 & $-$0{,}002 & $-$0{,}016 & -- & +0{,}005 \\
& {\scriptsize [$-$0{,}005; $-$0{,}002]} & {\scriptsize [$-$0{,}003; +0{,}000]} & {\scriptsize [$-$0{,}019; $-$0{,}014]} &  & {\scriptsize [+0{,}004; +0{,}005]} \\[2pt]
MUS & +0{,}024 & +0{,}004 & +0{,}017 & -- & +0{,}004 \\
& {\scriptsize [+0{,}022; +0{,}025]} & {\scriptsize [+0{,}002; +0{,}006]} & {\scriptsize [+0{,}015; +0{,}019]} &  & {\scriptsize [+0{,}004; +0{,}005]} \\[2pt]
PANDA & +0{,}007 & +0{,}001 & $-$0{,}022 & -- & +0{,}020 \\
& {\scriptsize [+0{,}005; +0{,}009]} & {\scriptsize [$-$0{,}001; +0{,}004]} & {\scriptsize [$-$0{,}024; $-$0{,}020]} &  & {\scriptsize [+0{,}019; +0{,}020]} \\[2pt]
Control 128\,px & +0{,}008 & +0{,}003 & +0{,}015 & -- & +0{,}000 \\
& {\scriptsize [+0{,}006; +0{,}010]} & {\scriptsize [$-$0{,}000; +0{,}006]} & {\scriptsize [+0{,}010; +0{,}019]} &  & {\scriptsize [$-$0{,}000; +0{,}001]} \\[2pt]
DTD & $-$0{,}003 & +0{,}017 & +0{,}015 & -- & +0{,}002 \\
& {\scriptsize [$-$0{,}004; $-$0{,}002]} & {\scriptsize [+0{,}015; +0{,}019]} & {\scriptsize [+0{,}013; +0{,}017]} &  & {\scriptsize [+0{,}001; +0{,}002]} \\[2pt]
\bottomrule
\end{tabular}
\end{table}


% Cifras clave (sensibilidad_pareada.csv):
% - C - B en TUM+NORM: KNN -0,016 [-0,018; -0,014]; MDS -0,019 [-0,021; -0,017];
%   RMDS +0,005 [+0,003; +0,006]; STOOD-X -0,041 [-0,043; -0,039].
% - Cambios en los dos sentidos: MDS en DTD +0,017; RMDS en LYM -0,016 y en PANDA -0,022.
% - ViM 256 - 171: de -0,001 (NORM) a +0,020 (PANDA); TUM+NORM +0,002.
\pendiente{comentario breve}

\section{Conclusión}\label{sec:apd-conclusion}

% [propuesta] Lo que se mantiene en los dos modelos:
% [propuesta] - MSP > MaxLogit > Energy en 9 de 9 conjuntos.
% [propuesta] - ReAct, ASH y DICE por debajo de MSP en 9 de 9.
% [propuesta] - LYM ~ TUM/NORM: AUROC(LYM) - AUROC(TUM+NORM) entre -0,003 y +0,021 en ConvNeXt (8 métodos).
%   El control interno apunta a las condiciones de adquisición en los dos modelos. Es la conclusión más
%   robusta del \cref{cap:deteccion-ood}.
% [propuesta] - Los métodos de distancia rinden más en PANDA que en el control (+0,156 a +0,207 en ConvNeXt).
% [propuesta] Lo que depende del modelo:
% [propuesta] - El mejor método: ViM y MDS en el Swin; KNN en ConvNeXt (0,972 en TUM+NORM).
% [propuesta] - MDS > KNN se invierte: KNN > MDS en 7 de 9 (IC sin el 0); en MUS y PANDA, el IC contiene el 0.
% [propuesta] - El fallo de RMDS en LYM es propio del Swin (ConvNeXt 0,970).
% [propuesta] - De los cuatro conjuntos en que STOOD-X era el primero con el Swin, sigue siéndolo en STR y MUS
%   y queda segundo, tras KNN, en NORM y DTD.
%   OJO: con ConvNeXt, STOOD-X también es el primero en PANDA (0,975) y en el control (0,782); «el primero
%   solo en STR y MUS» no es cierto sobre los nueve conjuntos.
% [propuesta] - MSP mucho mejor en ConvNeXt (+0,078 [+0,074; +0,082] en TUM+NORM): la ventaja de los métodos
%   basados en características sobre los logits es menor en este modelo (MSP 0,954 frente a MDS 0,953).
% [propuesta] - STOOD-X: la saturación es mucho menor en ConvNeXt (TUM+NORM 66,3 % -> 26,0 %; control
%   19,0 % -> 1,5 %) y sí distingue PANDA del control (0,859 frente a 0,475 en el Swin). Refuerza que la
%   coincidencia del Swin en PANDA y el control se debía a la saturación (\cref{sec:ood-panda}). La
%   saturación depende del espacio de características, no solo del método (en C, 9,9 % en TUM+NORM).
% [propuesta] Limitaciones:
% [propuesta] - Espacios no equivalentes: en C cambian algunos AUROC (STOOD-X -0,041 en TUM+NORM).
% [propuesta] - DICE colapsa en ConvNeXt (exactitud 0,2845 con el postprocesador de OpenOOD sin cambios).
% [propuesta] - CPU (ConvNeXt) frente a GPU (Swin).
% [propuesta] - Valores p sin corregir por comparaciones múltiples.
% [propuesta] - Los protocolos de entrenamiento de los dos modelos son distintos.
\pendiente{conclusión del autor}

\endinput
